\documentclass[10pt,a4paper]{article}

\usepackage[T1]{fontenc}
\usepackage[utf8]{inputenc}
\usepackage{lmodern}
\usepackage[english]{babel}
\usepackage{geometry}
\usepackage{microtype}
\usepackage{multicol}
\usepackage{enumitem}
\usepackage{titlesec}
\usepackage[hidelinks]{hyperref}
\usepackage{xcolor}
\usepackage{pgfgantt}
\usepackage{tabularx}
\usepackage{booktabs}

\definecolor{core}{RGB}{0,100,166}
\definecolor{opt}{RGB}{230,140,30}
\definecolor{deadline}{RGB}{190,30,45}
\newcommand{\legendbox}[1]{\tikz\fill[#1] (0,0) rectangle (0.3,0.2);}

\geometry{a4paper,margin=1.85cm}

\setlength{\parindent}{0pt}
\setlength{\parskip}{2.6pt}
\setlength{\columnsep}{0.75cm}

\setlist[itemize]{
  topsep=2pt,
  itemsep=0pt,
  parsep=0pt,
  leftmargin=1.4em
}

\titleformat{\section}
  {\normalsize\bfseries}{}{0pt}{}

\titlespacing*{\section}{0pt}{4pt}{2pt}

\pagestyle{empty}

\begin{document}


% TITLE
\begin{flushleft}

{\LARGE\bfseries
Project Proposal:
Tree Counts and Woody Species Diversity
from SWISSIMAGE and PlanetScope
\par}

\vspace{7pt}

{\small
Geospatial Analysis, Terrain \& Water (C) \quad | \quad
Jan Landolt, landoja2@students.zhaw.ch \quad | \quad
25. November 2026
}

\vspace{5pt}
\hrule

\end{flushleft}

\vspace{5pt}

% ABSTRACT (FULL WIDTH)
\section*{Abstract}

Hedgerows are key habitat structures in the Swiss
agricultural landscape, and their ecological value
depends strongly on woody species composition: the
Swiss quality level~II for hedges uses the number of
tree and shrub species per 10~m as a criterion. Field
assessment is time-consuming and rarely done at
landscape scale.

This project develops a workflow that provides, for
each hedge, (i)~the number of trees, estimated by
marker-controlled watershed segmentation of SWISSIMAGE
orthophotos (10~cm), and (ii)~the number of spectrally
distinct crown types, derived by colour clustering, as
a proxy for woody species richness. As an optional
extension, PlanetScope time series (3~m, near-daily)
could reveal seasonal colour trajectories (green-up,
flowering, autumn colouring) that separate species
looking alike in a single summer image.
\vspace{3pt}

% TWO-COLUMN BODY
\begin{multicols}{2}

\section*{Research Questions}

\emph{
(1) How well can individual trees per hedge be
delineated and counted from SWISSIMAGE using
marker-controlled watershed segmentation?
}

\emph{
(2) To what extent does the number of colour clusters
per hedge plausibly reflect its woody species
composition?
}

\emph{
(3, optional) Do PlanetScope phenological trajectories
improve the separation of species groups compared with
single-date colour?
}

\section*{Objectives}

\begin{itemize}

\item Delineate tree crowns and derive tree count and
density (trees per 100~m) per hedge.

\item Cluster crowns by colour and texture and derive
cluster number and spectral diversity per hedge.

\item Check both visually against SWISSIMAGE and, where
available, Kanton Zug hedge data for a subset of
hedges. Optionally test PlanetScope phenology.

\end{itemize}


\section*{Delimitation from Parent Project}

The parent project is my bachelor thesis
\emph{Automated detection and mapping of hedgerow
structures with LiDAR data -- potential and limits for
monitoring in Switzerland} (ZHAW/WSL, submission
January 2027). It derives hedge polygons and general structural
attributes for Kanton Zug from airborne LiDAR and gives
recommendations for the LABES monitoring.

The thesis asks \emph{where} hedges are and \emph{how
they are built}; this project asks \emph{what they
consist of}. It primary uses optical imagery, works at
crown level inside the given polygons and applies
methods that are not part of the thesis. The thesis polygons
serve as input only and they are neither re-derived nor evaluated here.

\vspace{2pt}
{\footnotesize
\begin{tabularx}{\columnwidth}{@{}>{\bfseries}p{1.25cm}
  >{\raggedright\arraybackslash}X
  >{\raggedright\arraybackslash}X@{}}
\toprule
 & Bachelor Thesis) & This project \\
\midrule
Question & Where are hedges, how are they built?
  & What do hedges consist of? \\
Data & Airborne LiDAR & SWISSIMAGE RGB, PlanetScope \\
Methods & Point-cloud metrics, classification of
  hedge vs.\ non-hedge
  & Watershed, colour clustering, phenology \\
\bottomrule
\end{tabularx}
}


\section*{Data and Methods}

The workflow will be implemented in Python or R
(scikit-image, scikit-learn, rasterio) and visualised with
QGIS. Data: SWISSIMAGE 10~cm (RGB), PlanetScope surface
reflectance (3~m, near-daily, incl. NIR),
hedge polygons (Landolt, 2027 [Bachelor-Thesis]) and, where
available, hedge data from Kanton Zug (e.g.\ species lists)
as visual reference.

Local maxima within a window matching the expected
crown diameter serve as markers for the watershed
segmentation. Segments below a minimum crown area are
removed; crowns are counted per hedge and 10~m section.

Second, colour (and texture features)
are extracted per crown and all data are grouped (k-means or
Gaussian mixture models), the number of clusters being
selected by silhouette score or BIC. The number of
clusters per hedge and a Shannon index of cluster
shares serve as spectral diversity measures.

Validation: for time reasons, validation is visual
only. Crown segments and colour clusters are overlaid
on SWISSIMAGE and, where available, compared with
Kanton Zug hedge data for a subset of hedges. Typical
errors (missed, merged or split crowns; clusters
mixing species) are described qualitatively. No field
survey and no statistical accuracy assessment are
carried out.


\section*{Optional Extension: Phenology}

Cloud-masked PlanetScope scenes (March to November)
are converted into smoothed vegetation index curves
per hedge section. Start of season, flowering and
autumn colouring could separate species with similar
summer colours, e.g.\ blackthorn (flowering before
leaf-out) and hawthorn (flowering later).


\section*{Expected Results and Limitations}

Expected results are a reproducible workflow, crown
maps and a visually checked table per hedge (tree
count, density, cluster number, spectral diversity).

Watershed segmentation of RGB imagery counts visible
crowns, not stems: overlapping crowns and dense shrubs
will be under-segmented, and trees and shrubs are hard
to separate without height data. Colour
clusters are not species, as illumination, shadow,
vitality and acquisition date affect colour.
PlanetScope pixels are coarse relative to hedge width,
causing mixed pixels with adjacent fields. The visual
validation gives no quantitative accuracy measure.

{\footnotesize
\textbf{References.}

Alonzo, M., et al. (2023).
\emph{Canopy composition drives variability in urban
growing season length more than the heat island
effect}.
Science of the Total Environment, 884, 163818.
\href{https://doi.org/10.1016/j.scitotenv.2023.163818}
{doi:10.1016/j.scitotenv.2023.163818}.

Ke, Y., \& Quackenbush, L.\,J. (2011).
\emph{A review of methods for automatic individual
tree-crown detection and delineation from passive
remote sensing}.
International Journal of Remote Sensing, 32(17),
4725--4747.
\href{https://doi.org/10.1080/01431161.2010.494184}
{doi:10.1080/01431161.2010.494184}.

Rocchini, D., et al. (2010).
\emph{Remotely sensed spectral heterogeneity as a
proxy of species diversity: recent advances and open
challenges}.
Ecological Informatics, 5(5), 318--329.
\href{https://doi.org/10.1016/j.ecoinf.2010.06.001}
{doi:10.1016/j.ecoinf.2010.06.001}.

Swiss Federal Council (2013).
\href{https://www.fedlex.admin.ch/eli/cc/2013/765/de}
{\emph{Direct Payments Ordinance (DZV, SR~910.13)}},
Annex~4.
}

\end{multicols}


% TIMELINE (FULL WIDTH)
\section*{Timeline}

Two variants: \textbf{(A)} core project only (Rsearch Questions~1--2)
with buffer, and \textbf{(B)} including the optional
phenology extension (Rsearch Questions~3). Variant~B is started only if
crown segmentation runs on all hedges by the go/no-go
date (11.12.); otherwise variant~A applies. Both account for the thesis
submission on the same day as the report (14.01.2027).

\vspace{7pt}

\vspace{4pt}
{\footnotesize
\legendbox{core}~Core \quad
\legendbox{opt}~Optional \quad
\legendbox{black!45}~Writing \quad
\legendbox{black!12}~Break / buffer \quad
\tikz\fill[deadline] (0,0) -- (0.12,0.12) -- (0.24,0)
  -- (0.12,-0.12) -- cycle;~Deadline
}

\ganttset{
  time slot format=isodate,
  x unit=0.19cm,
  y unit title=0.4cm,
  y unit chart=0.42cm,
  title label font=\scriptsize,
  title/.style={draw=black!30, fill=black!5},
  canvas/.style={draw=black!30},
  vgrid={*6{draw=none},*1{black!30,dotted}},
  bar height=0.6,
  bar top shift=0.2,
  bar/.style={draw=none, fill=core},
  bar label font=\footnotesize,
  milestone label font=\footnotesize,
  milestone/.style={shape=diamond, inner sep=0pt,
    draw=none, fill=black},
  milestone height=0.6,
  milestone top shift=0.2,
}
\tikzset{
  write/.style={draw=none, fill=black!45},
  brk/.style={draw=none, fill=black!12},
  optbar/.style={draw=none, fill=opt},
}
\newcommand{\gantthead}{%
  \gantttitle{Nov}{8}
  \gantttitle{December 2026}{31}
  \gantttitle{January 2027}{24} \\
  \gantttitle{CW48}{7} \gantttitle{CW49}{7}
  \gantttitle{CW50}{7} \gantttitle{CW51}{7}
  \gantttitle{CW52}{7} \gantttitle{CW53}{7}
  \gantttitle{CW1}{7} \gantttitle{CW2}{7}
  \gantttitle{CW3}{7} \\
}
\newcommand{\ms}[3][black]{%
  \ganttmilestone[milestone/.style={shape=diamond,
    inner sep=0pt, draw=none, fill=#1}]{#2}{#3}}

\vspace{6pt}
\textbf{Variant A -- core only (Rsearch Questions 1--2)}\par
\vspace{3pt}
\begin{ganttchart}{2026-11-23}{2027-01-24}
\gantthead
\ms{Proposal, coaching}{2026-11-25} \ms{}{2026-12-02} \\
\ganttbar{Data preparation}{2026-11-26}{2026-12-04} \\
\ganttbar{Crown segmentation (RQ 1)}{2026-12-03}{2026-12-12} \\
\ganttbar{Colour clustering (RQ 2)}{2026-12-09}{2026-12-17} \\
\ganttbar{Validation \& figures}{2026-12-14}{2026-12-19} \\
\ganttbar[bar/.style=write]{Report writing}{2026-12-16}{2026-12-19}
  \ganttbar[bar/.style=write]{}{2027-01-05}{2027-01-10} \\
\ganttbar[bar/.style=brk]{Christmas break / buffer}{2026-12-20}{2027-01-04}
  \ganttbar[bar/.style=brk]{}{2027-01-11}{2027-01-13} \\
\ms[deadline]{Report \& thesis due, exam}{2027-01-14} \ms{}{2027-01-18}
\end{ganttchart}

\vspace{6pt}
\textbf{Variant B -- with optional phenology (Rsearch Questions 1--3)}\par
\vspace{3pt}
\begin{ganttchart}{2026-11-23}{2027-01-24}
\gantthead
\ms{Proposal, coaching}{2026-11-25} \ms{}{2026-12-02} \\
\ganttbar{Data prep.\ incl.\ PlanetScope}{2026-11-26}{2026-12-04} \\
\ganttbar{Crown segmentation (RQ 1)}{2026-12-03}{2026-12-10} \\
\ganttbar{Colour clustering (RQ 2)}{2026-12-08}{2026-12-15} \\
\ms[opt]{Go/no-go optional part}{2026-12-11} \\
\ganttbar[bar/.style=optbar]{Phenology curves (RQ 3)}{2026-12-12}{2026-12-19} \\
\ganttbar[bar/.style=optbar]{Phenology vs.\ colour}{2027-01-05}{2027-01-07} \\
\ganttbar{Validation \& figures}{2026-12-15}{2026-12-19}
  \ganttbar{}{2027-01-06}{2027-01-08} \\
\ganttbar[bar/.style=write]{Report writing}{2027-01-07}{2027-01-12} \\
\ganttbar[bar/.style=brk]{Christmas break / buffer}{2026-12-20}{2027-01-04}
  \ganttbar[bar/.style=brk]{}{2027-01-13}{2027-01-13} \\
\ms[deadline]{Report \& thesis due}{2027-01-14} \ms{}{2027-01-18}
\end{ganttchart}


% RISK ANALYSIS (FULL WIDTH)
\newpage
\section*{Risk Analysis \& Plan B}

\textbf{SWOT.}
\vspace{2pt}

{\footnotesize
\setlist[itemize]{nosep, leftmargin=1.1em}
\begin{tabularx}{\textwidth}{@{}
  >{\raggedright\arraybackslash}X
  >{\raggedright\arraybackslash}X@{}}
\toprule
\textbf{Strengths}
\begin{itemize}
\item Hedge polygons and study area ready from the
  thesis.
\item SWISSIMAGE is free and nationwide: workflow is
  transferable to all of Switzerland.
\item Simple, well-documented methods (watershed,
  k-means/GMM) on an open Python stack.
\end{itemize}
&
\textbf{Weaknesses}
\begin{itemize}
\item Counts visible crowns, not stems.
\item Colour clusters are not always pecies.
\item Visual validation only, no quantitative
  accuracy.
\end{itemize}
\\ \midrule
\textbf{Opportunities}
\begin{itemize}
\item Composition indicator complements LiDAR structure
  for hedge quality (DZV quality level~II).
\item Phenology may separate early-flowering species
  (blackthorn vs.\ hawthorn).
\end{itemize}
&
\textbf{Threats}
\begin{itemize}
\item Cloud gaps in PlanetScope during key windows
  (spring flowering).
\item Mixed pixels: hedges are often narrower than
  2--3 PlanetScope pixels.
\item Illumination and date differences between
  SWISSIMAGE flight strips.
\end{itemize}
\\ \bottomrule
\end{tabularx}
}

\vspace{6pt}
\textbf{Risks and Plan B.}
\vspace{2pt}

{\footnotesize
\begin{tabularx}{\textwidth}{@{}
  >{\raggedright\arraybackslash}p{4.4cm}
  >{\raggedright\arraybackslash}p{1.6cm}
  >{\raggedright\arraybackslash}X@{}}
\toprule
\textbf{Risk} & \textbf{Likelihood} &
  \textbf{Mitigation / Plan B} \\
\midrule
Time conflict with thesis submission (14.01.2027)
  & High
  & Front-load analysis to December; start optional
    part only after go/no-go on 11.12.; fall back to
    variant~A. \\
Watershed over-/under-segments dense hedges
  & Medium
  & Tune smoothing and window size on one training
    hedge. Plan B: report crown density per 10~m
    section instead of single-tree counts. \\
Colour clusters do not reflect species richness
  & Medium
  & Report as a negative result. Plan B: test
    texture-only and CIELAB features, aggregate to
    species groups (e.g.\ conifer vs.\ broadleaf). \\
No Kanton Zug reference data for the study hedges
  & Medium
  & Plan B: visual check on SWISSIMAGE only; species
    interpretation of clusters stays qualitative. \\
PlanetScope cloud gaps or mixed pixels
  & High
  & Use only wide hedges and an inner buffer.
    Plan B: drop the optional part; core results
    (Rsearch Questions~1--2) are not affected. \\
Colour differences between SWISSIMAGE tiles
  & Medium
  & Normalise colour per hedge; compare hedges from the
    same flight date only. \\
\bottomrule
\end{tabularx}
}


% COACHING QUESTIONS (LAST)
\section*{Questions for the Coaching Session}

\begin{enumerate}[topsep=2pt, itemsep=1pt, leftmargin=1.4em]

\item \textbf{Delimitation:} Is it acceptable to use the
hedge polygons from my thesis as input? And could the
LiDAR canopy height model serve as Plan~B for crown
markers, or does that blur the line to the parent
project?

\item \textbf{Scope:} Is the core project (Rsearch Questions~1--2)
sufficient for the module, with phenology as an
optional extension?

\item \textbf{Validation:} Is a purely visual
validation (SWISSIMAGE, Kanton Zug data where
available) acceptable for the module, or is a small
quantitative check expected?

\item \textbf{Software:} Is a Python workflow
(scikit-image, scikit-learn) fine, or should the
analysis be implemented in R?

\item \textbf{Block:} Does the project fit Terrain \&
Water (raster analysis), or rather Spatial Ecology \&
Occurrence?

\end{enumerate}

\end{document}
